
%\needspace{4\baselineskip}
\color{uniblue}{\subsection{Защита по скорости изменения частоты (ЗИЧ)}}
\color{black}

\begin{enumerate}[label=\arabic{section}.\arabic{subsection}.\arabic*, labelsep=4pt, leftmargin=0pt, itemindent=57pt]

\item
Защита по скорости изменения частоты предназначена для фиксации резких изменений частоты в целях предупреждения развивающихся аварий в частотной системе. Данная функция основана на расчете скорости изменения частоты, функция не срабатывает при возникновении переходных нарушений в подаче напряжения. В устройстве защита по скорости изменения частот представлена независимыми функциями (ступенями) «ЗИЧ 1» и «ЗИЧ 2».

\item
Основные характеристики функций:
\begin{itemize}
    \item погрешность срабатывания ИО скорости изменения частоты должна быть не более 150 мГц;
    \item время срабатывания ИО скорости изменения частоты должно быть не более 200 мс.
\end{itemize}

\item
Функции содержат в себе ИО, реагирующий на скорость повышения частоты и скорость снижения частоты. Уставкой SGF2 предоставляется возможность выбора режима работы функций:
\begin{itemize}
    \item если уставка выставлена в положение «по dFdt>», то режим работы по скорости увеличения частоты;
    \item если уставка выставлена в положение «по dFdt<», то режим работы по скорости снижения частоты.
\end{itemize}

\item
Наличие выходного сигнала «ЗИЧ 1(2): ИО dF/dt» информирует о срабатывании ИО частоты по скорости изменения в составе функции.

\item
Срабатывание ЗИЧ происходит после отсчета выдержки времени, задаваемого уставкой T1.

\item
Если до окончания времени таймера выдержки времени на срабатывание скорость изменения частоты увеличивалась или снижалась в зависимости от положения регулируемой уставки SGF2, то активируется таймер задержки на возврат ИО Т2. Таймер задержки на возврат дает возможность таймеру выдержки времени на срабатывание продолжить работу, даже в случае, если скорость изменения частоты кратковременно становилась выше или ниже уставки.

\item
Блокировка функций по напряжению осуществляется снижением линейных напряжений ниже уставки, информируя выходным сигналом «ЗИЧ 1(2): Блок по U».

\item
Ввод функций в работу на этапе параметрирования устройства осуществляется своим программным переключателем SGF1 «Ввод\_функции». Информация о введенной функции отображается выходным сигналом «ЗИЧ 1(2): Ввод».

\item
Функции можгут быть выведены из работы оперативно путем активации сигнала «ОВ ЗИЧ 1(2)», что характеризуется активным состоянием сигнала «ЗИЧ 1(2): Оперативный вывод».

\item
При активном состоянии входного сигнала «ЗИЧ 1 на сигн», «ЗИЧ 2 на сигн» действие соответствующей функции ЗИЧ переводится на сигнал.

\item
Алгоритм работы ЗИЧ 1 выполнен в соответствии с рисунком \ref{zich:img1}. Алгоритм работы ЗИЧ 2 аналогичный. В таблице \ref{zich:tbl1} приведены параметры, необходимые для настройки ЗИЧ 1. В таблице \ref{zich1:tbl1} приведены параметры, необходимые для настройки ЗИЧ 2.

\vspace{1mm}
\begin{figure}[H]
\centering
\includegraphics[width=1\textwidth,height=1\textheight,keepaspectratio]{\fbpath/118. ЗИЧ/_latex/img/zich.pdf}
\captionof{figure}{Функциональная схема логики ЗИЧ 1}\label{zich:img1}
\end{figure}

\end{enumerate}
